\documentclass[11pt]{article}
\usepackage[a4paper,margin=25mm]{geometry}
\usepackage[T1]{fontenc}
\usepackage[utf8]{inputenc}
\usepackage{amsmath,amssymb,booktabs}
\usepackage{enumitem}
\usepackage[hidelinks]{hyperref}

\setlength{\parindent}{0pt}
\setlength{\parskip}{0.55em}
\setlist{nosep,leftmargin=*}
\emergencystretch=3em

\title{M13-C3 working fixture dossier:\\
binary deterministic AND collider\\
\texttt{m13\_bucket3\_binary\_collider\_and}}
\author{}
\date{Pre-run mathematical account (learning analysed separately)}

\begin{document}
\maketitle

\noindent\textbf{Status.} Active M1.3 Bucket~3 investigation; fixture certified;
sample \(n{=}30\), seed~\texttt{42} frozen; AAMAS+ECAI six-cell learning
\textbf{complete} (see \texttt{learning\_analysis.md} / \texttt{.tex});
\textbf{no Bucket~3 claim}.
Companion evidence record:
\texttt{docs/experiments/qualitative/M13-C3-binary-collider-and/experiment.md}.
This file remains the working \emph{pre-run} mathematical dossier, not
report-facing claim prose.

\section{Objects kept distinct}

\begin{enumerate}
  \item Generating DAG \(G\).
  \item Structural / conditional mechanisms.
  \item Exact observational population \(P\).
  \item Fixture-specific support, Markov, and ordinary-faithfulness analysis.
  \item Finite IID sample \(\mathcal{D}_{n{=}30}\).
  \item Later table-derived learner inputs (BK, \(E^+\), \(E^-\)).
  \item Later ABALearn outputs (delta / final framework).
  \item Evaluator-only mechanism reference and bounded interpretation.
\end{enumerate}
This pre-run dossier records items 1--5 and 8. Items 6--7 were produced later;
see \texttt{learning\_analysis.md} / \texttt{.tex}.

\section{Generating graph}

Internal learner-safe names: \(a,b,c\). Mathematical display: \(A,B,C\).

\[
G:\qquad A \rightarrow C \leftarrow B
\qquad\text{(no \(A\)--\(B\) edge).}
\]

Edges in the authoritative YAML:
\texttt{[a,c]}, \texttt{[b,c]}.
Unshielded collider at \(C\). Topological order \(a,b,c\).
Roles: \(A,B\) sources; \(C\) sink.

\section{Mechanisms}

Regime: \texttt{root\_stochastic\_deterministic\_nonroots} (schema version~2).

\begin{align*}
A &\sim \mathrm{Bernoulli}(4/5),
& P(A{=}0)&=1/5,& P(A{=}1)&=4/5,\\
B &\sim \mathrm{Bernoulli}(7/10),
& P(B{=}0)&=3/10,& P(B{=}1)&=7/10,\\
C &= A \land B
&&\text{(deterministic point-mass CPT).}
\end{align*}

Roots are mutually independent by fixture design
(\texttt{exogenous\_noise}: \texttt{mutually\_independent\_root\_exogenous\_noise}).
Randomness is confined to roots; \(C\) consumes no stochastic draw.

Truth table for \(C\mid(A,B)\):

\begin{center}
\begin{tabular}{@{}cccc@{}}
\toprule
\(A\) & \(B\) & \(C\) & \(P(A{=}a,B{=}b)\) \\
\midrule
0 & 0 & 0 & \(3/50\) \\
0 & 1 & 0 & \(7/50\) \\
1 & 0 & 0 & \(6/25\) \\
1 & 1 & 1 & \(14/25\) \\
\bottomrule
\end{tabular}
\end{center}

Both edges are \emph{active} in the local mechanism sense (certificate).
Active edges are not by themselves a faithfulness proof.

\section{Exact population}

Joint factorisation \(P(a,b,c)=P(a)P(b)P(c\mid a,b)\) yields four positive atoms
and four structural zeros (\texttt{population.csv}, \texttt{certificate.json}):

\begin{center}
\begin{tabular}{@{}lcc@{}}
\toprule
Assignment \((A,B,C)\) & Exact \(P\) & Decimal \\
\midrule
\((0,0,0)\) & \(3/50\) & \(0.06\) \\
\((0,1,0)\) & \(7/50\) & \(0.14\) \\
\((1,0,0)\) & \(6/25\) & \(0.24\) \\
\((1,1,1)\) & \(14/25\) & \(0.56\) \\
\midrule
structural zeros &
\((0,0,1),(0,1,1),(1,0,1),(1,1,0)\) & mass \(0\) \\
\bottomrule
\end{tabular}
\end{center}

Support size \(4\); joint state space size \(8\); \texttt{full\_support}: false;
minimum positive joint probability \(3/50\).

\section{Certificates}

Primary artefact:
\texttt{causal/outputs/causal\_fixtures/m13\_bucket3\_binary\_collider\_and/certificate.json}.
Scope: exact population; learner-visible flag \texttt{false}.

\subsection{Causal sufficiency}
Status: \texttt{declared\_satisfied\_by\_fixture\_design}.
Basis: mutually independent root exogenous noise declared in the YAML.
Not observationally testable from the joint alone; not inferred from CPTs.

\subsection{Causal Markov condition}
Status: \texttt{verified}.
Basis: joint derived as the product of validated local CPTs over \(G\),
including point-mass rows for \(C\). Graph-implied CI violations: \(0\).

\subsection{Ordinary faithfulness}
Status: \texttt{verified\_exactly}; \texttt{faithful}: true.
Basis: exact population CI set equals the graph \(d\)-separation set over the
exhaustive singleton-pair audit (six queries).
Explicitly \emph{not} implied by determinism or factorisation alone.

\subsection{Conditional independences}
Graph-implied and population independences coincide on the singleton audit:

\[
A \perp B
\qquad\text{(only population independence).}
\]

Dependence witnesses exist for \(A\not\!\perp B\mid C\), \(A\not\!\perp C\),
\(A\not\!\perp C\mid B\), \(B\not\!\perp C\), and \(B\not\!\perp C\mid A\),
as expected for this collider.

\subsection{MEC and ordinary CI CPDAG}
Markov equivalence class size \(1\); sole member is \(G\).
Ordinary graph-theoretic CI CPDAG:

\[
A \rightarrow C \leftarrow B
\qquad\text{(no undirected edges).}
\]

Certificate boundary note (verbatim intent): this is the ordinary CI-based CPDAG;
for deterministic variables, functional constraints may change what is identifiable
beyond ordinary CI; no extended recovery object is asserted here.

\section{Evaluator-only mechanism reference}

Artefact:
\texttt{mechanism\_reference.json}.
Encoding reference: \texttt{binary\_one\_vs\_zero\_exact\_value}
(positive state \(1\), negative state \(0\)).

\begin{center}
\begin{tabular}{@{}llp{0.48\textwidth}@{}}
\toprule
Target & Status & Canonical positive rule(s) \\
\midrule
\(a\) & \texttt{no\_observed\_parent\_deterministic\_rule} &
none (stochastic root) \\
\(b\) & \texttt{no\_observed\_parent\_deterministic\_rule} &
none (stochastic root) \\
\(c\) & \texttt{available} &
\texttt{c(A) :- a\_val\_1(A), b\_val\_1(A).} \\
\bottomrule
\end{tabular}
\end{center}

For \(C\), the formal, population-supported, and sample-observed positive rules
coincide on the string above. The artefact states that syntactic rule equality is
not an automatic causal-recovery verdict, and that predictive rules for root
targets do not constitute root-mechanism recovery.
\textbf{This reference must not enter learner BK or examples.}

\section{Frozen sample}

Approved baseline: \(n{=}30\), seed~\texttt{42}
(\texttt{samples/n30\_seed42.csv}).
Target-free IID observational rows; sampler
\texttt{row\_major\_ancestral\_inverse\_cdf}~v2; NumPy \texttt{PCG64}.
All \(30\) deterministic assignments verified.

Observed joint counts (all four population atoms appear):

\begin{center}
\begin{tabular}{@{}lc@{}}
\toprule
\((A,B,C)\) & Count \\
\midrule
\((1,1,1)\) & \(17\) \\
\((1,0,0)\) & \(7\) \\
\((0,1,0)\) & \(4\) \\
\((0,0,0)\) & \(2\) \\
\bottomrule
\end{tabular}
\end{center}

Empirical marginals: \(A{=}1\) in \(24/30\); \(B{=}1\) in \(21/30\); \(C{=}1\) in \(17/30\).

\section{Implied target counts under the pre-run binary policy}

At the pre-run stage, the planned target-wise policy used target value \(1\) as
\(E^+\) and value \(0\) as \(E^-\). It was later executed as recorded in
\texttt{learning\_analysis.md}. Counts implied by the frozen sample marginals:

\begin{center}
\begin{tabular}{@{}lrr@{}}
\toprule
Target & \(\lvert E^+\rvert\) & \(\lvert E^-\rvert\) \\
\midrule
\(a\) & \(24\) & \(6\) \\
\(b\) & \(21\) & \(9\) \\
\(c\) & \(17\) & \(13\) \\
\bottomrule
\end{tabular}
\end{center}

These are \textbf{not} generated learner example arrays and have
\textbf{not} been supplied to ABALearn. No BK, delta, or solution artefacts
exist yet for this fixture under
\texttt{causal/outputs/aba\_learning/targetwise/}.

\section{Provenance (selected hashes)}

From \texttt{fixture\_manifest.json}:

\begin{itemize}
  \item Source UTF-8 bytes:
    \texttt{sha256:4ec5bb5726f4359a6b347e519ebcdbf1ef79a945ccd54a40d601dad240d824f8}
  \item Canonical parsed YAML document:
    \texttt{sha256:f061e29abc8ba8bc1fcd8630ef1f79e38c8171b36b6ee09afb5f88e94eac4298}
  \item Semantic model:
    \texttt{sha256:e76c1ff5b849cf6f11ec161e4dbd23573fb39891d846596b155bf047fd687731}
  \item Sample CSV:
    \texttt{sha256:1edae465d0f6b6b662b3c4ed4863af21c1d9b0c1cf428774731dd3f43282a8ca}
\end{itemize}

\texttt{model.bif} is derived interoperability only (floating-point probabilities);
it is not authoritative.

\section{Current status and next step}

AAMAS and ECAI six-cell target-wise learning on the frozen sample is complete.
Learning-stage narrative:
\texttt{learning\_analysis.tex} / \texttt{learning\_analysis.md}
(investigation hub: \texttt{experiment.md}).
\textbf{No Bucket~3 claim.}
Follow-up probes H1--H3 are complete / analysed; remaining probe H4 is recorded
in \texttt{future\_probes.tex} / \texttt{future\_probes.md}. None is
part of the closed H0 six-cell analysis.

\end{document}
